%% The MIT License (MIT)
%%
%% Copyright (c) 2015 Daniil Belyakov
%%
%% Permission is hereby granted, free of charge, to any person obtaining a copy
%% of this software and associated documentation files (the "Software"), to deal
%% in the Software without restriction, including without limitation the rights
%% to use, copy, modify, merge, publish, distribute, sublicense, and/or sell
%% copies of the Software, and to permit persons to whom the Software is
%% furnished to do so, subject to the following conditions:
%%
%% The above copyright notice and this permission notice shall be included in all
%% copies or substantial portions of the Software.
%%
%% THE SOFTWARE IS PROVIDED "AS IS", WITHOUT WARRANTY OF ANY KIND, EXPRESS OR
%% IMPLIED, INCLUDING BUT NOT LIMITED TO THE WARRANTIES OF MERCHANTABILITY,
%% FITNESS FOR A PARTICULAR PURPOSE AND NONINFRINGEMENT. IN NO EVENT SHALL THE
%% AUTHORS OR COPYRIGHT HOLDERS BE LIABLE FOR ANY CLAIM, DAMAGES OR OTHER
%% LIABILITY, WHETHER IN AN ACTION OF CONTRACT, TORT OR OTHERWISE, ARISING FROM,
%% OUT OF OR IN CONNECTION WITH THE SOFTWARE OR THE USE OR OTHER DEALINGS IN THE
%% SOFTWARE.

% The font could be set to Windows-specific Calibri by using the 'calibri' option
\documentclass[]{mcdowellcv}

% For mathematical symbols
\usepackage{amsmath}

% Set applicant's personal data for header
\name{Ashwin Iyer}
\address{Boston, MA}
\availability{Availability: May -- Aug 2027}
\contacts{(682) 239-9481 \linebreak iyer.ashw@northeastern.edu \linebreak \href{https://ashwiniyer.com}{ashwiniyer.com}}

\usepackage[hidelinks]{hyperref}

\begin{document}

	% Print the header
	\makeheader
	\vspace{-4mm}
	\begin{cvsection}{Education}
		\begin{cvsubsection}{ Northeastern University}{Boston, MA}{Expected May 2028}
		\textbf{Candidate for Bachelor of Science in Computer Science and Business Administration \hfill{GPA: 3.8}}
		\textbf{Honors/Activities:} NU Systematic Alpha, Dean's List
		\vspace{-2mm}
			
		\textbf{Relevant Coursework:} Discrete Structures, Introduction to Databases, Program Design \& Implementation, Business Statistics, Financial Management
		\end{cvsubsection}
	\end{cvsection}
	% Print the content
\vspace{-4mm}

\begin{cvsection}{Technical Skills}
	\begin{cvsubsection}{}{}{}

\textbf{Languages:} Python, C++, Java, Rust, JavaScript, TypeScript, SQL, Kotlin
\vspace{-2mm}

\textbf{Frameworks \& Libraries:} React, Next.js, Redux, Electron, TensorFlow, Keras, Pandas, NumPy
\vspace{-2mm}

\textbf{Developer Tools:} Git, AWS (EC2), PostgreSQL, Azure DevOps, IntelliJ, Eclipse, PyCharm, Xcode
	\end{cvsubsection}
\end{cvsection}
\vspace{-4mm}

			\begin{cvsection}{Work Experience}
		\begin{cvsubsection}{ Global Risk \& Analytics Co-op}{}{December 2025 -- June 2026}
			\textit{Wellington Management | Boston, MA}
			\begin{itemize}
				\item Built Python risk management tools that apply proprietary \textbf{factor risk models} to compute risk metrics across equity and alternative asset portfolios.
				\item Developed a \textbf{Model Context Protocol (MCP) server} that lets internal LLM workflows query risk metrics, run \textbf{stress-test scenarios}, and flag portfolios breaching mandate limits, adding a compliance guardrail to the risk platform.
				\item Produced \textbf{top-down portfolio risk reports} surfacing concentration risk across multi-asset portfolios.
			\end{itemize}
		\end{cvsubsection}
\vspace{-2mm}

		\begin{cvsubsection}{ Software Engineering Intern}{}{May 2025 -- August 2025}
			\textit{Zeal IT Consultants | Dallas, TX}
			\begin{itemize}
				\item Developed the frontend for Trinity Industries' Asset Management System using React and Next.js.
				\item Increased UI development velocity by \textbf{10+ story points per sprint} and grew team delivery capacity by \textbf{300\%} within one release cycle, accelerating the project timeline by 4 weeks.
				\item Cut page load times by \textbf{94\%} by migrating state management from MobX to Redux and implementing server-side rendering.
			\end{itemize}
		\end{cvsubsection}
	\end{cvsection}
	\vspace{-4mm}

	\begin{cvsection}{Projects}
		\begin{cvsubsection}{}{}{}	
                \textbf{Event-Driven Implied Volatility Mispricing} | \textit{Python, Polygon.io} | \underline{\href{https://ashwiniyer.com/projects/NUSA\%20Spring\%2026\%20RS\%20-\%20Equity\%20Vol.\%20Project.pdf}{Report}} \hfill{January 2026 -- March 2026}
				\begin{itemize}
					\item Built a Python backtesting framework extending Gao, Xing \& Zhang (2018) to identify event-driven implied volatility (IV) mispricings across \textbf{4} event types on \textbf{541} historical S\&P 500 constituents and \textbf{15} sector ETFs (2016--2025), processing \textbf{80,000+} option observations.
					\item Generated sector-level long/short straddle signals achieving a \textbf{63\%} win rate and \textbf{3.62\%} average return per trade out-of-sample on Technology earnings, and a \textbf{74\%} win rate with \textbf{10.81\%} ROI on FOMC events in 2025 testing.
				\end{itemize}
				\smallskip
                \textbf{Prediction Market Trading} | \textit{Rust, AWS} | \underline{\href{https://kalshi.com/ideas/profiles/whao}{Portfolio}} \hfill{December 2025 -- Present}
				\begin{itemize}
					\item Implemented a mathematical model in Rust that prices a prediction market in real time, hosted on AWS EC2 for low-latency API access.
					\item Identified a market edge and scaled the strategy to a \textbf{net adjusted Sharpe ratio of 1.2} over two months, with \textbf{40\%} total returns and a maximum drawdown of \textbf{10\%}.
				\end{itemize}
				\smallskip
					\textbf{PM-Trading Desk} | \textit{Python, WebSockets} | \underline{\href{https://github.com/Ashwin-Iyer1/pm-tradingdesk}{Github}} \hfill{September 2025 -- Present}
				\begin{itemize}
					\item Developed a prediction market trading application using WebSockets for real-time data access, processing \textbf{500+} market updates per second with \textbf{sub-50ms} order execution latency.
					\item Implemented a \textbf{smart order router} spanning Polymarket and Kalshi, achieving \textbf{2--5\%} average price improvement per trade through \textbf{cross-exchange execution}.
				\end{itemize}
				\smallskip
					\textbf{PaveGuard} | \textit{React, Python, YOLO} | \underline{\href{https://github.com/Ashwin-Iyer1/Pothole-detector}{Github}} \hfill{October 2023}
				\begin{itemize}
					\item Trained a YOLO image recognition model on \textbf{8,000+} images to classify potholes and road fractures across \textbf{5} damage types with \textbf{71\%} accuracy.
					\item Built a full-stack crowdsourced reporting platform with a React frontend and Python backend, winning the \textbf{top prize} out of \textbf{40+} teams at the AI for All hackathon at UT Dallas.
				\end{itemize}
		\end{cvsubsection}
	\end{cvsection}
\vspace{-4mm}
	\begin{cvsection}{Interests}
		\begin{cvsubsection}{}{}{}
				Hackathons, Reading, Rubik's Cube, Chess, Poker, Baseball, Blogging, Football, Working Out, Watches, Shoes
		\end{cvsubsection}
	\end{cvsection}

	
\end{document}

